\باب{مخلوط اعداد}\شناخت{ضمیمہ_ج}
مخلوط اعداد سے مراد \عددی{a+ib} روپ کا اظہار ہے جہاں \عددی{a} اور \عددی{b} حقیقی اعداد ہیں اور \عددی{i} جذر \عددی{\sqrt{-1}} کی علامت ہے۔حقیقی اعداد کا نظام تخلیق دینے میں  جدوجہد درکار تھی جس  کے مختلف مراحل پر   اس ضمیمے میں غور کیا جائے گا۔ یہ جاننے کے بعد مخلوط اعداد کا نظام عجیب نظر  نہیں آئے گا۔

\جزوحصہء{حقیقی اعداد کا  ارتقا}
اعداد ی نظام  کے ارتقا میں پہلا قدم گنتی کے اعداد \عددی{1,2,3,\cdots} کی پہچان تھی جنہیں ہم اب\اصطلاح{قدرتی اعداد }\فرہنگ{اعداد!قدرتی}\حاشیہب{natural numbers}\فرہنگ{numbers!natural}یا \اصطلاح{مثبت عدد صحیح }\فرہنگ{اعداد!مثبت عدد صحیح}\حاشیہب{positive integers}\فرہنگ{integers!positive} کہتے ہیں۔ نظام کے اندر رہتے ہوئے ان اعداد کو استعمال کر کے چند بنیادی  حسابی  عملیات ممکن ہیں۔یعنی مثبت اعداد صحیح کا نظام جمع اور ضرب کے عمل  کے لحاظ سے\اصطلاح{بند }\فرہنگ{بند}\حاشیہب{closed}\فرہنگ{closed} نظام ہے۔ اس سے ہمارا مراد یہ ہے کہ اگر \عددی{m} اور \عددی{n} مثبت عدد صحیح ہوں تب  درج ذیل بھی مثبت عدد  صحیح   ہوں گے۔
\begin{align}\label{مساوات_ج_عدد_صحیح_بند_الف}
m+n=p\quad  \text{اور}\quad mn=q  
\end{align}
مساوات \حوالہ{مساوات_ج_عدد_صحیح_بند_الف} کی دونوں مساواتوں  میں بائیں ہاتھ دو عدد صحیح جانتے ہوئے ہم مطابقتی دائیں ہاتھ کے  عدد صحیح دریافت کر سکتے ہیں۔ بعض اوقات  مثبت عدد صحیح \عددی{m} اور \عددی{p} جانتے ہوئے ہم  ایسا مثبت عدد صحیح \عددی{n} تلاش کر سکتے ہیں کہ \عددی{m+n=p} ہو۔ مثال کے طور ہم  \عددی{3+n=7} کو حل کر کے  ہم  مثبت عدد صحیح  \عددی{n} تلاش کر سکتے ہیں۔ اس کے برعکس \عددی{7+n=3}   کو حل کر کے ہم ایسا \عددی{n} دریافت نہیں کر سکتے جو مثبت عدد صحیح  ہو۔ اس مثال سے واضح ہے کہ مثبت عدد صحیح نظام زیادہ کارآمد نہیں اور اس کو وسیع بنانے کی ضرورت ہے۔

منفی اعداد صحیح اور عدد  صفر کی ایجاد نے  \عددی{7+n=3} طرز کی مساوات حل کرنا ممکن بنایا۔  ایسی تہذیب میں، جو تمام
\اصطلاح{عدد صحیح}\فرہنگ{عدد صحیح}\حاشیہب{integers}\فرہنگ{integers}
\begin{align}\label{مساوات_ضمیمہ_ج_عدد_صحیح}
\cdots,-3,-2,-1,0,1,2,3,\cdots
\end{align}
 جانتی ہو ، ایک پڑھا لکھا شخص  مساوات \عددی{m+n=p} میں دو عدد صحیح جانتے ہوئے  نامعلوم عدد صحیح تلاش کر سکتا ہے۔

فرض کریں یہ پڑھا لکھا شخص   مساوات \حوالہ{مساوات_ضمیمہ_ج_عدد_صحیح} میں دیے گئے اعداد استعمال کر کے ضرب کرنا بھی جانتا ہے۔اگر  \عددی{m} اور \عددی{q} معلوم ہو، وہ دیکھے گا کہ بعض اوقات   \عددی{n} تلاش کرنا ممکن ہوتا ہے اور بعض اوقات ناممکن۔اگر یہ شخص قابل اور ذہین ہو، ممکن ہے وہ اعداد کے نظام کو مزید وسعت دے کر      عدد صحیح \عددی{m} اور \عددی{n} کی مرتب جوڑیاں  \عددی{m\!/\!n}،  جس کو ہم کسر  کہتے ہیں،  متعارف کرے۔عدد صفر  خصوصی خواص رکھتا ہے جو کچھ عرصہ کے لئے مسئلے پیدا کر سکتا ہے، لیکن وہ جلد اس نتیجے پر پہنچ پائیں گے کہ \عددی{m\!/\!n} طرز کی  تمام  جوڑیاں کارآمد ہیں تاہم    جن کے نسب نما میں صفر  ہو انہیں  عددی نظام میں شامل نہیں  کیا جائے۔ یہ 
مجموع ، جو \اصطلاح{ناطق اعداد  }\فرہنگ{اعداد!ناطق}\حاشیہب{rational numbers}\فرہنگ{numbers!rational} کہلاتا ہے ،  نظام میں کسی بھی دو اعداد  پر حساب کی     درج ذیل \اصطلاح{ناطق  عملیات}\فرہنگ{ناطق!عملیات}\حاشیہب{rational operations}\فرہنگ{rational operations}    ممکن بناتا ہے، تاہم \ترچھا{صفر سے تقسیم کرنا ممکن نہیں}۔
\begin{multicols}{2}
\begin{enumerate}[1.]
\item
ا) جمع\\
ب) منفی
\item
ا) ضرب\\
ب) تقسیم
\end{enumerate}
\end{multicols}

اکائی مربع کا ہندسہ (شکل \حوالہء{2}) اور مسئلہ فیثاغورث سے ہم جانتے ہیں کہ وہ  ہندسی  بناؤٹ کی مدد سے لمبائی کی کسی بنیادی اکائی کے لحاظ سے   \عددی{\sqrt{2}}  اکائی لمبائی  کی لکیر کھینچ سکتے تھے۔ یوں ہم کہہ سکتے ہیں کہ ہندسی  بناؤٹ سے وہ درج ذیل مساوات حل کر سکتے تھے۔
\[x^2=2\]
تاہم  ان کو معلوم ہوا کہ   \عددی{\sqrt{2}} لمبائی کی لکیر اور \عددی{1} لمبائی کی لکیر    متبائن  مقادیر ہیں۔ یعنی   \عددی{\sqrt{2}/1}   کو کسی دوسری زیادہ بنیادی اکائی لمبائی    کے دو عدد صحیح   مضرب کی نسبت  نہیں لکھا جا سکتا۔یوں ایک پڑھا لکھا شخص مساوات \عددی{x^2=2} کا حل ناطق عدد کی صورت میں دریافت نہیں کر سکتا تھا۔

 ایسا کوئی  ناطق عدد نہیں  پایا جاتا جس کا  مربع  \عددی{2} ہو۔ یہ جاننے کے لئے کہ ایسا کیوں کر ہے، فرض کریں  کہ ایسا  ناطق عدد پایا جاتا ہے۔یوں ہم ایسے  دو عدد صحیح \عددی{p} اور \عددی{q}  دریافت کر سکیں گے جن  کا \عددی{1} کے علاوہ کوئی مشترک جزو ضربی  نہ ہو اور  جو درج ذیل کو مطمئن کرتے ہوں۔
\begin{align}
p^2=2q^2
\end{align}
چونکہ \عددی{p} اور \عددی{q} عدد صحیح ہیں، \عددی{p}  کو  جفت ہونا ہو گا ورنہ اس کا اپنے ساتھ حاصل ضرب طاق ہو گا۔ علامتوں میں \عددی{p=2p_1} ہو گا جہاں \عددی{p_1} عدد صحیح ہے۔یوں \عددی{2p_1^2=q^2}  لکھا جا سکتا ہے جس کے تحت \عددی{q} کو جفت ہونا ہو گا مثلاً \عددی{q=2q_1^2}  جہاں \عددی{q_1} عدد صحیح ہے۔اس طرح \عددی{2} کو \عددی{p} اور \عددی{q} دونوں کا جزو ضربی ہونا ہو گا،  لیکن ہم نے \عددی{p} اور \عددی{q} کا انتخاب اس طرح کیا  تھا کہ \عددی{1} کے علاوہ ان کا کوئی مشترک جزو ضربی نہ ہو۔ لہٰذا ایسا کوئی ناطق عدد نہیں پایا جاتا جس کا مربع \عددی{2} ہو۔

اگرچہ  پڑھا لکھا شخص مساوات  \عددی{x^2=2}  کو حل نہیں کر سکتا تاہم وہ ناطق اعداد   کی ترتیب 
\begin{align}\label{مساوات_ضمیمہ_خیالی_چار}
\frac{1}{1}, \frac{7}{5}, \frac{41}{29}, \frac{239}{169}, \cdots
\end{align}
وضع کر سکتا تھا جن کے مربع درج ذیل ترتیب دیتے ہیں
\begin{align}\label{مساوات_ضمیمہ_خیالی_پانچ}
\frac{1}{1}, \frac{49}{25}, \frac{1681}{841}, \frac{57121}{28561}, \cdots
\end{align}
جس کی حد  \عددی{2} پر مرکوز ہے۔اس حقیقت  کی وجہ سے    انہیں ناطق اعداد کی ترتیب کی حد  کا تصور واضح ہوا۔ اگر ہم  اس حقیقت کو  مان لیں کہ اوپر سے محدود  بڑھتی  ترتیب   ہر صورت  ایک حد پر مرکوز ہو گی   اور  ہم دیکھیں کہ   (\حوالہ{مساوات_ضمیمہ_خیالی_چار})   میں دی گئی  ترتیب  یہ خواص رکھتی ہے تب ہم  چاہیں گے کہ یہ حد \عددی{L} ہو۔ ترتیب \حوالہ{مساوات_ضمیمہ_خیالی_پانچ}  کو دیکھ کر اس کا یہ مطلب بھی  ہے کہ \عددی{L^2=2} ہو گا لہٰذا \عددی{L}    ناطق عدد نہیں ہو سکتا۔  ناطق  اعداد کے ساتھ   ناطق اعداد کی تمام بڑھتی  محدود ترتیبات کی حدیں  شامل کرنے سے ہمیں تمام   حقیقی اعداد کا نظام ملتا ہے۔یہاں  لفظ \ترچھا{حقیقی}  اس نظام  کو کسی بھی لحاظ سے ریاضی کے دیگر  نظام سے زیادہ یا کم حقیقی نہیں بناتا۔ یہ محض ایک نام ہے جو اس نظام کو دیا گیا ہے۔

\جزوحصہء{مخلوط اعداد}
ہم نے  حقیقی  اعداد کا نظام  بنتے دیکھا ۔درحقیقت اس نظام  کے ارتقا میں  تین مراحل  دیکھے گئے۔
\begin{enumerate}[1.]
\item
پہلے عددی نظام  کی ایجاد: \quad
تمام عدد صحیح  کا  سلسلہ  سادہ گنتی سے حاصل کیا۔
\item
دوسرے  عددی نظام کی ایجاد:\quad
ناطق اعداد \عددی{m\!/!n } کا سلسلہ   اعداد صحیح  سے  حاصل کیا گیا۔
\item
تیسرے عددی نظام کی ایجاد:\quad
حقیقی اعداد \عددی{x} کا سلسلہ  ناطق اعداد اسے  حاصل کیا گیا۔
\end{enumerate}

عددی نظام  کی ایجادات  درجائی ہیں ۔ ہر نئے   نظام  میں گزشتہ نظام  شامل ہیں۔ہر نیا نظام زیادہ وسیع ہے اور اس نظام  میں  رہتے ہوئے  مزید  ریاضیاتی عمل ممکن ہیں۔

\begin{enumerate}[1.]
\item
تمام  اعداد صحیح  کے نظام میں درج ذیل روپ کی تمام مساوات حل کی جا سکتی ہیں جہاں \عددی{a} عدد صحیح ہے۔
\begin{align}\label{مساوات_ضمیمہ_مخلوط_چھ}
x+a=0
\end{align}
\item
تمام ناطق اعداد کے نظام میں درج ذیل روپ کی مساوات حل کی جا سکتی ہیں جہاں \عددی{a} اور \عددی{b} ناطق اعداد ہیں اور \عددی{a\ne 0} ہے۔
\begin{align}\label{مساوات_ضمیمہ_مخلوط_سات}
ax+b=0
\end{align}
\item
تمام حقیقی اعداد کے نظام میں مساوات \حوالہ{مساوات_ضمیمہ_مخلوط_چھ} اور مساوات \حوالہ{مساوات_ضمیمہ_مخلوط_سات} کے علاوہ  تمام دو درجی مساوات
\begin{align}\label{مساوات_ضمیمہ_مخلوط_آٹھ}
ax^2+bx+c=0
\end{align}
حل کرنا ممکن ہے، جہاں \عددی{a\ne 0}  اور \عددی{b^2-4ac\ge 0} ہیں۔ 
\end{enumerate}

آپ غالباً درج ذیل کلیہ جانتے ہیں جو مساوات \حوالہ{مساوات_ضمیمہ_مخلوط_آٹھ} کا حل دیتا ہے۔
\begin{align}\label{مساوات_ضمیمہ_مخلوط_نو}
x=\frac{-b\pm\sqrt{b^2-4ac}}{2a}
\end{align}
آپ یقیناً یہ بھی جانتے ہوں گے جب\اصطلاح{مفرق}\فرہنگ{مفرق}\حاشیہب{discriminant}\فرہنگ{discriminant}، \عددی{d=b^2-4ac}،  منفی ہو ، مساوات \حوالہ{مساوات_ضمیمہ_مخلوط_نو} کا حل  کسی بھی  مذکورہ بالا اعدادی نظام میں نہیں پایا جاتا۔ درحقیقت، مذکورہ بالا ایجاد شدہ اعداد نظام کے اندر رہتے ہوئے،  درج ذیل نہایت سادہ دو درجی  مساوات کا حل ممکن نہیں۔
\[x^2+1=0\]

اس طرح ہم چوتھا  عددی نظام ایجاد کرتے ہیں، جو  تمام مخلوط اعداد \عددی{a+ib} پر مشتمل ہے۔ہم علامت \عددی{i} کو رد کر کے علامت \عددی{(a,b)} استعمال کر سکتے ہیں، جہاں  ہم عدد صحیح  \عددی{a} اور \عددی{b}  کی جوڑی کی بات کریں گے۔الجبرائی عمل میں \عددی{a} اور \عددی{b}   کے ساتھ مختلف  برتاؤ کیا جاتا ہے لہٰذا  قوسین میں ان کی ترتیب  تبدیل نہیں  کی جا سکتی۔یوں ہم کہہ سکتے ہیں کہ مخلوط عددی نظام، اعداد صحیح کی  تمام مرتب جوڑیوں \عددی{(a,b)} پر مشتمل ہے  اور  جوڑیوں کی مساویت، مجموعہ، فرق، ضرب، وغیرہ درج ذیل قواعد کے تحت   ہیں۔  درج ذیل بحث میں  \عددی{(a,b)} اور \عددی{a+ib} علامتیت دونوں استعمال کی گئی ہے۔ ہم \عددی{a} کو مخلوط عدد \عددی{(a,b)} کا \اصطلاح{حقیقی جزو }\فرہنگ{حقیقی جزو}\حاشیہب{real part}\فرہنگ{real part} اور \عددی{b} کو اس کا\اصطلاح{خیالی جزو }\فرہنگ{خیالی جزو}\حاشیہب{imaginary part}\فرہنگ{imaginary part} کہتے ہیں۔

ہم درج ذیل متعارف کرتے ہیں۔

\ترچھا{مساویت}\\
\begin{minipage}[t]{0.45\textwidth}
دو مخلوط اعداد \عددی{a+ib} اور \عددی{c+id} صرف اور صرف  اس صورت ایک دوسرے کے  برابر: \\
{\centerline{\(a+ib=c+id\)}}\\
ہوں گے جب  \عددی{a=c} اور \عددی{b=d} ہو۔
\end{minipage}\hfill
\begin{minipage}[t]{0.45\textwidth}
دو مخلوط اعداد \عددی{(a,b)} اور \عددی{(c,d)}   صرف اور صرف اس صورت ایک دوسرے کے  برابر ہوں گے جب \عددی{a=c} اور \عددی{b=d} ہو۔
\end{minipage}


\ترچھا{جمع}\\
\begin{minipage}[t]{0.45\textwidth}
\((a+ib)+(c+id)\\=(a+c)+i(b+d)\)
\end{minipage}\hfill
\begin{minipage}[t]{0.45\textwidth}
دو مخلوط اعداد \عددی{(a,b)} اور \عددی{(c,d)} کا مجموعہ  مخلوط عدد \عددی{{(a+c,b+d)}}  ہو گا۔
\end{minipage}

\ترچھا{ضرب}\\
\begin{minipage}[t]{0.45\textwidth}
\((a+ib)(c+id)\\
=(ac-bd)+i(ad+bc)\)
\end{minipage}\hfill
\begin{minipage}[t]{0.45\textwidth}
دو مخلوط اعداد \عددی{(a,b)} اور \عددی{(c,d)} کا  حاصل ضرب  مخلوط عدد \عددی{{(ac-bd,ad+bc)}}  ہو گا۔
\end{minipage}

ان تمام مخلوط اعداد \عددی{(a,b)} کا سلسلہ جن میں دوسرا عدد \عددی{b} صفر ہو حقیقی اعداد \عددی{a} کے سلسلہ کے تمام خواص رکھتا ہے۔ مثال کے طور پر \عددی{(a,0)} اور \عددی{(b,0)}  کا مجموعہ اور حاصل ضرب درج ذیل ہوں گے
\begin{align*}
(a,0)+(c,0)&=(a+c,0)\\
(a,0)\cdot (c,0)&=(ac,0)
\end{align*}
جو اسی قسم کے اعداد ہیں اور ان کا خیالی جزو صفر ہے۔اسی طرح  \قول{حقیقی عدد} \عددی{(a,0) }اور مخلوط عدد \عددی{(c,d)} کو آپس میں ضرب کرنے سے درج ذیل حاصل ہو گا۔
\[(a,0)\cdot (c,d)=(ac,ad)=a(c,d)\]
بالخصوص،مخلوط عددی نظام میں  مخلوط عدد \عددی{(0,0)}   صفر کا کردار اور مخلوط عدد \عددی{(1,0)}   اکائی کا کردار ادا کرتا ہے۔

جوڑی عدد \عددی{(0,1)}، جس کا حقیقی جزو صفر  اور خیالی جزو ایک کے برابر ہے، کی خاصیت یہ ہے کہ اس کے مربع:
\[(0,1)\cdot(0,1)=(-1,0)\]
میں حقیقی جزو  منفی ایک اور خیالی جزو صفر کے برابر ہے۔یوں، مخلوط اعداد \عددی{(a,b)} کے نظام میں  ایسا عدد \عددی{x=(0,1)} پایا جاتا ہے جس کے  مربع  کو اکائی \عددی{(1,0)}کے ساتھ جمع کرنے سے    صفر\عددی{(0,0)}   حاصل ہوتا ہے (ذیل دیکھیں)۔
\[(0,1)^2+(1,0)=(0,0)\]
یوں اس نئے عددی نظام میں درج ذیل مساوات  کا حل \عددی{x=(0,1)} ہے۔
\[x^2+1=0\]

آپ غالباً  \عددی{(a,b)} کے لحاظ سے \عددی{a+ib}  کی  علامت سے زیادہ واقف ہیں۔ مرتب جوڑیوں کے  ریاضی قواعد کے تحت درج ذیل لکھا جا سکتا ہے
\[(a,b)=(a,0)+(0,b)=a(1,0)+b(0,1)\]
اور چونکہ \عددی{(1,0)}  کا رویہ اکائی کی طرح اور \عددی{(0,1)}  کا رویہ منفی ایک کے جذر کی طرح  ہے لہٰذا \عددی{(a,b)} کی جگہ \عددی{a+ib}  لکھنے میں جھجک محسوس نہ  کریں، جس میں \عددی{b} کے ساتھ چسپاں \عددی{i}  یقینی بناتا ہے کہ خیالی جزو کی درست نشاندہی ہو۔ ہم جب چاہیں \عددی{a+ib} کی جگہ \عددی{(a,b)} یا \عددی{(a,b)} کی جگہ \عددی{a+ib} لکھ سکتے ہیں۔ مخلوط عددی نظام \عددی{(a,b)}  میں رائج  الجبرائی قواعد جاننے کے بعد    یاد رہے کہ علامت  \عددی{(0,1)=i} کسی بھی لحاظ سے علامت \عددی{(1,0)} سے کم حقیقی نہیں ہے۔پس ایک کو حقیقی اور دوسرے کو خیالی نام  دیا  گیا ہے۔

مخلوط اعداد کے کسی بھی  ناطق  ملاپ  کو واحد ایک مخلوط عدد میں  سمیٹنے کے لئے بنیادی الجبرا کے قواعد استعمال کریں اور جہاں \عددی{i^2} ہو وہاں   \عددی{-1} لکھیں۔یاد رہے کہ  مخلوط عدد \عددی{(0,0)=0+i0} سے تقسیم کی اجازت نہیں۔ اگر \عددی{a+ib\ne 0} ہو تب تقسیم کا عمل درج ذیل ہو گا۔
\[\frac{c+id}{a+ib}=\frac{(c+id)(a-ib)}{(a+ib)(a-ib)}=\frac{(ac+bd)+i(ad-bc)}{a^2+b^2}\]
نتیجہ مخلوط عدد \عددی{x+iy} ہے جہاں \عددی{x} اور \عددی{y} درج ذیل ہیں
\[x=\frac{ac+bd}{a^2+b^2},\quad y=\frac{ad-bc}{a^2+b^2}\]
جہاں \عددی{a+ib=(a,b)\ne (0,0)} کی وجہ سے \عددی{a^2+b^2\ne 0} ہو گا۔

نسب نما    کو   ضارب \عددی{a-ib}   سے ضرب دے کر نسب  نما میں   \عددی{i} سے چھٹکارہ  حاصل کیا گیا۔ اس ضارب کو  \عددی{a+ib} کا \اصطلاح{مخلوط جوڑی دار}\فرہنگ{مخلوط جوڑی دار}\حاشیہب{complex conjugate}\فرہنگ{complex conjugate} کہتے ہیں۔   روایتی طور پر  مخلوط عدد \عددی{z} کے مخلوط جوڑی دار کو  \عددی{\bar{z}}  سے ظاہر کیا جاتا ہے۔ یوں درج ذیل ہو گا۔
\[z=a+ib,\quad \bar{z}=a-ib\]
کسر \عددی{(c+id)\!/\!(a+ib)} کے  شمارکنندہ اور نسب نما کو  مخلوط جوڑی دار سے ضرب دینے سے نسب نما  ہر صورت حقیقی عدد میں تبدیل ہو گا۔

\ابتدا{مثال}
\begin{enumerate}[a.]
\item\hfill$\begin{aligned}[t]
(2+3i)+(6-2i)=(2+6)+(3-2)i=8+i
\end{aligned}$\hfill\mbox{}
\item\hfill$\begin{aligned}[t]
(2+3i)-(6-2i)=(2-6)+(3-(-2))i=-4+5i
\end{aligned}$\hfill\mbox{}
\item\hfill$\begin{aligned}[t]
(2+3i)(6-2i)&=(2)(6)+(2)(-2i)+(3i)(6)+(3i)(-2i)\\
&=12-4i+18i-6i^2=12+14i+6=18+14i
\end{aligned}$\hfill\mbox{}

\item\hfill$\begin{aligned}[t]
\frac{2+3i}{6-2i}&=\left(\frac{2+3i}{6-2i}\right)\left(\frac{6+2i}{6+2i}\right)\\
&=\frac{12+4i+18i+6i^2}{36+12i-12i-4i^2}\\
&=\frac{6+22i}{40}=\frac{3}{20}+\frac{11}{20}i
\end{aligned}$\hfill\mbox{}
\end{enumerate}
\انتہا{مثال}


\جزوحصہء{اغگن خاکے}
مخلوط عدد \عددی{z=x+iy}  کا ہندسی اظہار  درج ذیل  دو طریقوں سے کیا جاتا ہے۔
\begin{enumerate}[a.]
\item
اس کو  مستوی \عددی{xy} میں  نقطہ \عددی{P(x,y)}  سے ظاہر کیا جاتا ہے۔
\item
اس کو مستوی \عددی{xy}  میں مبدا سے نقطہ \عددی{P} تک سمتیہ \عددی{\koverrightharp{OP}}سے ظاہر کیا جاتا ہے۔
\end{enumerate}
